\section{Complete-source coefficient transport on a fixed wall}
\label{sec:coefficient-transport}

The root velocity equations introduce inverse internal gaps. Passing to the
complete prepared polynomial removes those denominators from the transport
equation. The source information is still required, but it can be acquired on
one exterior wall. We give the finite theorem with its source cost explicit.

Fix a circle $\Gamma=\{|z-z_0|=R\}$, positively oriented, and a real time
interval $[t,T]$. Suppose $H$ is holomorphic on a neighborhood of the closed
spatial disk and the time interval, satisfies
\begin{equation}\label{eq:heat-fixed-wall}
 \partial_tH=-\partial_z^2H,
 \qquad H(u,w)\ne0\quad(u\in[t,T],\ w\in\Gamma),
\end{equation}
and has $N$ zeros inside $\Gamma$, counted with multiplicity. Write
$H=UP$, where $P$ is the complete monic polynomial of those zeros and $U$ is
holomorphic and nonzero in the disk. The center and radius remain fixed.
For a polynomial $f$, set
\[
 Q_f(w,z)=\frac{f(w)-f(z)}{w-z},\qquad
 Z_H=\frac{H_{zz}}H,\qquad
 \mathcal A_Hf=\frac1{2\pi i}\int_\Gamma Z_H(w)Q_f(w,z)\,dw.
\]
The divided difference $Q_f$ is a polynomial, including when $w=z$.

\begin{proposition}[Complete-source equation]\label{prop:source-coeff}
The prepared polynomial satisfies
\begin{equation}\label{eq:source-coeff}
 P_t=-\mathcal A_HP=-P_{zz}-2\mathcal T_P(bP_z),
 \qquad b=U_z/U,
\end{equation}
where
\[
 \mathcal T_Pg(z)=\frac1{2\pi i}\int_\Gamma
                   \frac{g(w)}{P(w)}Q_P(w,z)\,dw.
\]
These identities remain valid at internal multiple zeros.
\end{proposition}

\begin{proof}
The operator $\mathcal T_P$ gives the degree less than $N$ analytic remainder
modulo $P$; in particular it fixes every polynomial of degree less than $N$
and annihilates $P$ times a function holomorphic in the disk. These statements
follow directly from Cauchy's formula and do not require simple roots.
Since $H_t/H=P_t/P+U_t/U$, multiplying by $Q_P$ and integrating gives
$P_t=\int_\Gamma(H_t/H)Q_P/(2\pi i)$. The heat equation proves the first
identity. Finally
$H_{zz}/U=P_{zz}+2bP_z+(U_{zz}/U)P$, and analytic remainder gives the second.
\end{proof}

Put $\zeta=z-z_0$, and let $S$ be the plain coefficient shift:
$S\zeta^j=\zeta^{j-1}$ for $j\ge1$, and $S1=0$. It is not
$d/dz$. Define the complete source moments
\[
 M_k(t)=\frac1{2\pi i}\int_\Gamma(w-z_0)^k Z_H(t,w)\,dw.
\]
Expanding a divided monomial gives, on polynomials of degree at most $n$,
\begin{equation}\label{eq:source-shift}
 \mathcal A_H=\sum_{j=1}^{n}M_{j-1}S^j.
\end{equation}
Thus all these operators commute, including operators constructed from
different source fields. They lower degree and are nilpotent.

\begin{theorem}[Finite coefficient propagator]\label{thm:nilpotent-propagator}
Let $m_j(t,s)=\int_t^sM_{j-1}(u)\,du$. The backward propagator of
$E'=-\mathcal A_HE$ on degree at most $n$ is
\begin{equation}\label{eq:nilpotent-propagator}
 \mathcal U_H(t,s)=\exp\!\left(\sum_{j=1}^n m_j(t,s)S^j\right).
\end{equation}
The exponential terminates after degree $n$. For the fixed coefficient norm
$\|f\|_R=\sum_j|f_j|R^j$, its exact induced norm is
\begin{equation}\label{eq:propagator-norm}
 B_n(t,s)=\sum_{k=0}^n
 \left|[v^k]\exp\!\left(\sum_{j=1}^n\frac{m_j(t,s)}{R^j}v^j\right)\right|.
\end{equation}
If $V=\int_t^s\sup_\Gamma|Z_H(u,w)|\,du$, then
\begin{equation}\label{eq:propagator-majorant}
 B_n(t,s)\le\sum_{\ell=0}^n\binom n\ell\frac{V^\ell}{\ell!}.
\end{equation}
\end{theorem}

\begin{proof}
Equation~\eqref{eq:source-shift} makes the time generators commuting
polynomials in $S$, so differentiating the finite exponential proves
\eqref{eq:nilpotent-propagator}. If $\mathcal U=\sum b_kS^k$, the weighted
column sum of its $j$th column is $\sum_{k=0}^j|b_k|R^{-k}$, maximal at
$j=n$. This proves the exact norm formula. The wall estimate gives
$|m_j|/R^j\le V$. Majorize the generating series by
$\exp(Vv/(1-v))$ through degree $n$. Summing those coefficients and using
$\sum_{k=\ell}^n\binom{k-1}{\ell-1}=\binom n\ell$ proves
\eqref{eq:propagator-majorant}.
\end{proof}

The estimate is polynomial in acquired integrated source sizes at fixed
degree. It is not asserted to be polynomial uniformly in height, population
size, or a late-time parameter. For a difference of monic degree-$N$
polynomials, use $n=N-1$ in the error propagator, while propagating the
polynomial itself uses $n=N$.

\subsection{Reference errors and what must be acquired}

Let $G$ be a reference holomorphic field with the same zero count and with
$G(u,w)\ne0$ for every $u\in[t,T]$ and $w\in\Gamma$, with complete monic prepared polynomial $P_*$. It need not be an exact
heat solution. Define
\[
 d_*=\partial_tP_*+\mathcal A_GP_*,\qquad E=P-P_*,\qquad
 r=(\mathcal A_H-\mathcal A_G)P_*+d_*.
\]
Then $E'=-\mathcal A_HE-r$, and hence
\begin{equation}\label{eq:paid-duhamel}
 E(t)=\mathcal U_H(t,T)E(T)+\int_t^T\mathcal U_H(t,s)r(s)\,ds.
\end{equation}
The complete reference defect is also a wall quantity:
\[
 d_*(z)=\frac1{2\pi i}\int_\Gamma
                   \frac{G_t+G_{zz}}G(w)Q_{P_*}(w,z)\,dw.
\]
If $\eta_Z\ge\sup_\Gamma|Z_H-Z_G|$ and
$W_R(P_*)=\sum_jj|p_{*,j}|R^j$, a closed finite error estimate is
\begin{equation}\label{eq:paid-error}
 \|E(t)\|_R\le B_{N-1}(t,T)\|E(T)\|_R+
 \int_t^TB_{N-1}(t,s)
       \bigl(\eta_Z(s)W_R(P_*(s))+\|d_*(s)\|_R\bigr)\,ds.
\end{equation}
A directly evaluated correlated residual can replace this absolute majorant.
Its size must be acquired independently; an unknown coefficient difference
inside a forcing bound is not free data.

For example, acquired field and second-derivative errors $e_0,e_2$, reference
wall floor $m_*>e_0$, and reference second-derivative upper bound $M_{2,*}$ give
\begin{equation}\label{eq:source-quotient-budget}
 \eta_Z\le\frac{e_2}{m_*-e_0}
              +\frac{M_{2,*}e_0}{m_*(m_*-e_0)}.
\end{equation}
An analytic field-error bound on a spatial collar of width $d$ gives
$e_2\le2e_{\rm collar}/d^2$. This Cauchy step concerns the field error,
not a logarithmic ratio continued through different interior zero sets.

For exact heat fields the integrated moments have the endpoint representation
\begin{equation}\label{eq:source-endpoint-log}
 m_j(t,T)=-\frac1{2\pi i}\int_\Gamma(w-z_0)^{j-1}
              \log\frac{H(T,w)}{H(t,w)}\,dw,
\end{equation}
using the logarithm lifted continuously from time $t$. Thus equally informed
endpoint wall data already determine the homogeneous propagator. The
commutation theorem does not confer an intrinsic advantage on a curved path;
a gain would have to come from improved source estimates or conditioning along
that path, with its entire error paid.

\begin{remark}
The fixed-source generator is different from free polynomial heat evolution
on arbitrary error vectors. For $H=z^3-6tz$, one finds
$\mathcal A_H=6S^2$ on degree at most three, while
$\partial_z^2(z^2)=2$. Comparing two independently evolving polynomial heat
fields includes the source-quotient difference in
\eqref{eq:paid-duhamel}. Conversely, suppressing the external unit in
\eqref{eq:source-coeff} changes the complete problem. Note~12 contains the
free-heat comparison, complex-path formulation and collision controls.
\end{remark}

\section{A complete reference at microscopic real heat times}
\label{sec:microscopic-reference}

We now specify a theta source input rather than assume a small wall error.
Throughout,
\[
 H_t(z)=\int_0^\infty e^{tu^2}\Phi(u)\cos(zu)\,du,
 \qquad H_0(z)=\tfrac18\xi(\tfrac12+\tfrac{iz}2).
\]
The effective approximation of Polymath~\cite{polymath} applies for real
$0<t\le1/2$, $x\ge200$, $0\le y\le1$. In this section $z=x+iy$ and
\[
 N=\left\lfloor\sqrt{x/(4\pi)+t/16}\right\rfloor.
\]
We work on a chart where $N$ stays fixed. Changes of $N$ require explicit
patching; a jumping truncation is not a holomorphic reference.

Define, with principal logarithms away from their cuts,
\[
 M_0(s)=\frac{s(s-1)}{16}\pi^{-s/2}\sqrt{2\pi}
 \exp\!\left[(s/2-1/2)\operatorname{Log}(s/2)-s/2\right],
\quad
 \alpha(s)=\frac1{2s}+\frac1{s-1}+\frac12\operatorname{Log}\frac{s}{2\pi}.
\]
Here $\alpha=M_0'/M_0$. Put $s=(1-iz)/2$ and
\begin{equation}\label{eq:complete-finite-reference}
 G_{t,N}(z)=\sum_{\sigma\in\{s,1-s\}}\sum_{n=1}^N
 M_0(\sigma)n^{-\sigma}
             \exp\!\left[\frac t4(\alpha(\sigma)-\log n)^2\right],
 \qquad B_t(z)=M_0(s)e^{t\alpha(s)^2/4}.
\end{equation}
This rewrites the complete two-sum reference without complex conjugates.
In particular $G$ is locally holomorphic in $t,z$, although it is not an exact
heat solution.

\subsection{Explicit field and heat-defect budgets}

For completeness, a finite coarsening of the imported error bounds is as
follows. Set $X=x/(4\pi)$, $L=\log X$, and
\[
 p=\operatorname{Re}(s+t\alpha(s)/2),\quad
 \rho=e^{0.02y}X^{-y/2},\quad \nu=\frac{ty}{2(x-6)},
\]
\[
 S_0=\sum_{n=1}^N(1+\rho N^\nu n^y)e^{t\log^2n/4}n^{-p},
 \quad b_0=\frac{t^2L^2/16+0.626}{x-6.66},\quad
 E_{AB}=S_0(e^{b_0}-1).
\]
The exact value of $p$ avoids relying on slightly different printed
lower-bound corrections for its real part. Theorem~1.3 of
\cite{polymath} gives $|H_t-G_{t,N}|/|B_t|\le E_{AB}+E_{C,0}$, with
\[
 E_{C,0}=X^{-(1+y)/4}
 \exp\!\left[-\frac{tL^2}{16}
 +\frac{1.24(3^y+3^{-y})}{N-0.125}
 +\frac{3|L+i\pi/2|+10.44}{x-12}\right].
\]
For $tL^2\le K$ and $y=0$, $E_{AB}=O_K(X^{-3/4})$, whereas
$E_{C,0}=O_K(X^{-1/4})$. Indeed the retained weights have modulus at most
a constant times $n^{-1/2}$, and their relative error bracket is $O_K(1/x)$.
These are usable absolute errors even when $tL^2\to0$; the late-clock
assumption $x\ge e^{C/t}$ is not needed. That assumption is needed for the
much shorter two-carrier asymptotic, whose omitted absolute sum is not small
on the microscopic clock.

For a term $g=M_0(\sigma)n^{-\sigma}e^{t\beta^2/4}$, with
$\beta=\alpha(\sigma)-\log n$, direct differentiation gives
\begin{equation}\label{eq:finite-reference-defect}
 (\partial_t+\partial_z^2)g=-\frac g4
 \left[\alpha'+t\beta^2\alpha'
 +\frac t2\bigl((\alpha')^2+\beta\alpha''\bigr)
 +\frac{t^2}{4}\beta^2(\alpha')^2\right].
\end{equation}
Here both $\sigma=s$ and $\sigma=1-s$ satisfy
$\partial_z^2=-\partial_\sigma^2/4$. Taking absolute values in this finite
sum is an explicit defect budget. Since $\beta=O(\log x)$,
$\alpha'=O(x^{-1})$, $\alpha''=O(x^{-2})$, it proves
\begin{equation}\label{eq:uncorrected-defect-order}
 |G_t+G_{zz}|/|B_t|=O_K(X^{-3/4})\qquad(y=0,\ tL^2\le K).
\end{equation}
No source quotient has been divided by an unacquired wall floor in this step.

\subsection{An analytic continuation of the explicit correction}

Corollary~6.4 and Proposition~6.6(vi) of \cite{polymath} improve the field
remainder by retaining a correction $C_t$. Define
\[
 v_0=x/2+\pi t/8,\quad p_N(v)=1+2N-\sqrt{2v/\pi},
\]
\[
 U(v)=\exp\!\left[-i\left(\frac v2\log\frac v{2\pi}-\frac v2-\frac\pi8\right)\right],
 \quad
 C_0(p)=\frac{e^{\pi i(p^2/2+3/8)}-i\sqrt2\cos(\pi p/2)}{2\cos(\pi p)}.
\]
The singularities at $p=\pm1/2$ are removable, with value $(1-i)/4$.
The displayed source correction is
\[
 C_t(x+iy)=2(-1)^Ne^{-\pi iy/8+t\pi^2/64}
       \operatorname{Re}\{M_0(iv_0)C_0(p_N(v_0))U(v_0)e^{\pi i/8}\}.
\]
It gives
\begin{equation}\label{eq:corrected-field-remainder}
 \frac{|H_t-(G_{t,N}-C_t)|}{|B_t|}\le E_{AB}+E_C,
\end{equation}
where
\[
 E_C=X^{-(1+y)/4}
 e^{-tL^2/16+(3|L+i\pi/2|+3.58)/(x-8.52)}
 \left[\frac{1.24(3^y+3^{-y})}{N-0.125}+\frac{6.92}{x-12}\right].
\]
The extra $1/N$ makes this remainder $O_K(X^{-3/4})$ on the real axis.
However, the displayed $C_t$ is generally not holomorphic in $z$; its real-part
operation cannot be differentiated as an analytic reference.

There is an explicit analytic replacement. Let
$C_0^\#(p)=\overline{C_0(\bar p)}$,
$C_R=(C_0+C_0^\#)/2$, $C_I=(C_0-C_0^\#)/(2i)$, and
\[
 V_N(v)=v^{3/2}C_R(p_N(v))-v^{1/2}C_I(p_N(v)),\qquad
 k_N=-(-1)^N\frac{\sqrt\pi}{4}.
\]
The functions $C_R,C_I$ are holomorphic continuations of the real and imaginary
parts on the real axis, not real-part operations on complex arguments. Set
\begin{equation}\label{eq:analytic-correction}
 \widetilde C_{t,N}(z)=k_Ne^{-\pi z/8-\pi^2t/64}
                      V_N(z/2+\pi t/8),\qquad
 \widehat G=G-\widetilde C.
\end{equation}
Use principal square roots on $\operatorname{Re}v>0$.

\begin{proposition}[Corrected analytic reference]\label{prop:corrected-reference}
The correction \eqref{eq:analytic-correction} agrees with $C_t$ at $y=0$.
With
\[
 Q(t,z)=\frac{\sqrt\pi}{4}\frac{e^{-\pi x/8-\pi^2t/64}}{|B_t(z)|},
\]
it satisfies the explicit bounds
\begin{align}
 \frac{|H_t-\widehat G|}{|B_t|}
 &\le E_{AB}+E_C+Q(t,z)\frac{|y|}{2}
       \sup_{0\le\theta\le1}|V_N'(v_0+i\theta y/2)|,
                                                   \label{eq:analytic-correction-error}\\
 (\partial_t+\partial_z^2)\widetilde C
 &=\frac{k_N}{4}e^{-\pi z/8-\pi^2t/64}
                         V_N''(z/2+\pi t/8).       \label{eq:analytic-correction-defect}
\end{align}
For fixed $K$ and $y_0<1$, its normalized field allowance is
$O_{K,y_0}(X^{-(3+y)/4})$ uniformly for $0\le y\le y_0$ and
$0<tL^2\le K$. Its correction defect is
$O_{K,y_0}(X^{-(5+y)/4})$. In particular both the field error and total heat
defect of $\widehat G$ are $O_K(X^{-3/4})$ at $y=0$.
\end{proposition}

\begin{proof}
For positive real $v$, substituting
$\operatorname{Log}(iv/2)=\log(v/2)+i\pi/2$ gives the exact cancellation
\[
 M_0(iv)U(v)e^{\pi i/8}
       =-\frac{\sqrt\pi}{8}\sqrt v\,(v+i)e^{-\pi v/4}.
\]
Consequently
$C_t(x+iy)=k_Ne^{-\pi(x+iy)/8-\pi^2t/64}V_N(v_0)$.
The fundamental theorem of calculus from $v_0$ to $v_0+iy/2$ proves
\eqref{eq:analytic-correction-error}. In the heat defect, the coefficients of
$V_N$ and $V_N'$ cancel because the argument derivatives are $1/2$ and
$\pi/8$, while the exponential logarithmic derivatives are $-\pi/8$ and
$-\pi^2/64$. This proves \eqref{eq:analytic-correction-defect}.

For completeness the derivative sizes have explicit, zero-independent
constants. Along the comparison segment,
\[
 |p_N(v)-p_N(v_0)|\le\frac{|y|}{2\sqrt{2\pi v_0}}<\frac18,
\]
so $|p_N(v)|<9/8$. On this disk both $C_R^{(j)}$ and $C_I^{(j)}$ are bounded
by $K_j=j!8^jM_*$ for $j=0,1,2$, where
\[
 M_*=\frac{e^{25\pi/32}+\sqrt2\cosh(5\pi/8)}{2\sinh(\pi/8)}.
\]
Indeed $|\cos\pi p|\ge\sinh(\pi/8)$ on $|p|=5/4$: split into
$|\operatorname{Im}p|\ge1/8$ and its complement, where
$1<|\operatorname{Re}p|\le5/4$. Bound the numerator on that circle, remove
$\pm1/2$, and apply Cauchy's estimate with radius $1/8$.
Since $p_N'=-1/\sqrt{2\pi v}$ and $p_N''=-p_N'/(2v)$, differentiation gives
$V_N'=O(|v|)$ and $V_N''=O(|v|^{1/2})$ with these explicit constants.
Finally $Q=O_{K,y_0}(X^{-(7+y)/4})$ follows directly from the displayed
Stirling factor. Combine these estimates with
\eqref{eq:finite-reference-defect}. Note~14 records the full derivative
formulas and convenient finite majorants.
\end{proof}

\subsection{The remaining wall and time gates}

These are absolute source allowances. To use the coefficient theorem, acquire
a wall floor $m_f>e_0$ for $f=\widehat G/B$, where $e_0$ bounds the
right-hand side of \eqref{eq:analytic-correction-error}. Then
$|H/\widehat G-1|\le e_0/m_f$. For a common analytic normalization write
$e=H/B-f$ and $\lambda=B'/B$. The complete quotient cancellation is
\[
 \frac{H_{zz}}H-\frac{\widehat G_{zz}}{\widehat G}
 =\frac{fe''-ef''+2\lambda(fe'-ef')}{f(f+e)}.
\]
Thus an acquired spatial collar supplies derivative errors without introducing
internal root gaps. The local reference derivative ratios still require
bounds. A collar crossing the real axis must convert the reflected error to
one analytic carrier before applying Cauchy's estimate.

All imported field errors here concern real $t$. The explicit references
extend analytically in complex time on their charts, but that does not extend
the error for $H$ with the same bound. A complex-time route needs a separately
paid source remainder. Finally a fixed-$N$ chart is not the same as a fixed
number of roots in a selected block: a prepared population count, common wall,
and the forest denominator remain separate acquisition obligations. The finite
certificate in Note~13 pays those obligations in its original bounded box.
No estimate in this section supplies the missing unbounded-height small-gap
upper bound or a global simplicity conclusion.
